\section{An exact test for one-sided hidden \texorpdfstring{$F_3$}{F3} violation}\label{sec:orbit}
The centre bound is a universal obstruction. Using the entire ellipsoid gives
a necessary and sufficient test for whether any Bob-only filter can reveal
$F_3>1$. This directly addresses the asymmetric filtering problem raised
in Eqs.~(9)--(10) and the conclusion of the 2024 paper~\cite{ducuara}.

Fix canonical Bloch data $(\cc,0,S)$ and put
\begin{equation}\label{eq:orbitdefs}
\tau=\norm S_F^2,\qquad \mathbf k=S^{\mathsf T}\cc,\qquad
K=(\tau-c^2)\Id-S^{\mathsf T}S.
\end{equation}
Every positive invertible Bob filter, after normalization, is represented by
$F=\sqrt{\Id+\bb\cdot\bs}$, $\norm\bb<1$. Polar decomposition shows that
post-filter unitaries do not enlarge the achievable $F_3$ values. From
\eqref{eq:filteridentity} and $\aa=\cc+S\bb$, the exact squared value is
\begin{equation}\label{eq:orbitquadratic}
f(\bb)=\tau+2\mathbf k\cdot\bb-\bb^{\mathsf T}K\bb.
\end{equation}
On the unit sphere its continuous boundary value is
\begin{equation}\label{eq:orbitboundary}
f(\bb)=\norm{\cc+S\bb}^2\leq1.
\end{equation}
Boundary parameters represent singular filters producing product states.

\begin{theorem}[One-sided violation criterion and exact optimum]\label{thm:orbit}
Some invertible Bob filter produces $F_3>1$ if and only if
\begin{equation}\label{eq:orbitcriterion}
K\succ0,\qquad \norm{K^{-1}\mathbf k}^2<1,\qquad
V:=\tau+\mathbf k^{\mathsf T}K^{-1}\mathbf k>1.
\end{equation}
When these conditions hold, the maximum squared value is $V$, attained at
the unique positive-filter parameter $\bb_*=K^{-1}\mathbf k$.
Singular or indefinite $K$ cannot permit an $F_3$ violation.
\end{theorem}
\begin{proof}
A value greater than one in the open ball gives, by continuity, a maximum
greater than one on the compact closed ball. Equation~\eqref{eq:orbitboundary}
forces it to be interior. First and second derivative conditions give
$K\bb_*=\mathbf k$ and $K\succeq0$.
If $K$ had a nonzero null vector $\mathbf v$, then
$\mathbf k\cdot\mathbf v=\bb_*^{\mathsf T}K\mathbf v=0$.
The value of $f$ would be constant on the line $\bb_*+s\mathbf v$, which
meets the unit sphere, contradicting \eqref{eq:orbitboundary}.
Thus $K\succ0$. Completing the square gives
\[
f(\bb)=V-(\bb-\bb_*)^{\mathsf T}K(\bb-\bb_*).
\]
This proves the necessary conditions, value, and uniqueness. Conversely,
\eqref{eq:orbitcriterion} places this stationary point in the physical
open ball with value greater than one; the same identity proves optimality.
\end{proof}

For an original state with full-rank Bob marginal, whitening followed by this
filter is one Bob filter, scaled to a valid Kraus contraction. Conversely
every Bob filter is represented in this canonical orbit up to a Bob unitary.
A pure Bob marginal is a product case and never violates. Thus both
possible marginal ranks and every global rank are covered.

\subsection{Ellipsoid-only and exact polynomial forms}
Put $L=(\tau-c^2)\Id-Q$, where $Q=SS^{\mathsf T}$. The matrices $K,L$ have
the same eigenvalues, including when $S$ is singular. Intertwining yields
\begin{equation}\label{eq:geometriccriterion}
L\succ0,\qquad \cc^{\mathsf T}QL^{-2}\cc<1,\qquad
\tau+\cc^{\mathsf T}QL^{-1}\cc>1.
\end{equation}
This equivalent test depends only on the centre and ellipsoid matrix.
For rational data, let $\ell=\det L$ and $R=\operatorname{adj}L$.
After certifying $L\succ0$ by principal minors, so $\ell>0$, the other tests are
\[
\cc^{\mathsf T}QR^2\cc<\ell^2,\qquad
(\tau-1)\ell+\cc^{\mathsf T}QR\cc>0.
\]
These are finite polynomial inequalities and require no numerical inverse.

The unit-ball condition cannot be dropped. The positive definite canonical state
\[
\sigma=\frac1{40}\begin{pmatrix}19&0&0&5\\0&7&1&0\\0&1&1&0\\5&0&0&13\end{pmatrix}
\]
has two positive $2\times2$ blocks with determinants $222/1600$ and $6/1600$.
Its data are $\cc=(0,0,3/10)$, $S=\operatorname{diag}(3/10,-1/5,3/5)$,
$\tau=49/100$, $K=\operatorname{diag}(31,36,4)/100\succ0$.
The unconstrained vertex is $(0,0,9/2)$ with value $13/10>1$, outside the
physical ball. The actual supremum is $81/100<1$, from
\[
f(x,y,z)=\frac{49}{100}+\frac9{25}z-\frac{31}{100}x^2
-\frac9{25}y^2-\frac1{25}z^2.
\]
It increases with $z$ on $[-1,1]$, and its transverse terms are nonpositive.

The criterion gives another proof of the threshold. For $c>0$ and
$q=\nn^{\mathsf T}Q\nn$, $\nn=\cc/c$, some direction perpendicular to $\nn$
has Rayleigh quotient at least $(\tau-q)/2$. Hence
\[
\lambda_{\min}(L)\leq\frac{\tau+q-2c^2}{2}
\leq\frac{1-3c^2+3q}{2}\leq2-3c.
\]
For $c\geq2/3$, the necessary positive-definiteness condition fails.

\subsection{The entire orbit supremum, including singular quadratic cases}
For completeness, the exact scalar trust-region characterization is
\begin{equation}\label{eq:trustregion}
\sup_{\norm\bb<1}f(\bb)=\tau+
\min_{\substack{\lambda\geq0,\ K+\lambda\Id\succeq0\\
\mathbf k\in\operatorname{range}(K+\lambda\Id)}}
\left[\lambda+\mathbf k^{\mathsf T}(K+\lambda\Id)^\dagger\mathbf k\right].
\end{equation}
Here $\dagger$ is the Moore--Penrose inverse. For the upper bound, put
$B=K+\lambda\Id$ and complete the square:
\[
f(\bb)=\tau+\lambda\norm\bb^2+\mathbf k^{\mathsf T}B^\dagger\mathbf k
-(\bb-B^\dagger\mathbf k)^{\mathsf T}B(\bb-B^\dagger\mathbf k).
\]
For equality, use $\lambda=0$ if $K\succeq0$,
$\mathbf k\in\operatorname{range}K$, and $\norm{K^\dagger\mathbf k}\leq1$.
Otherwise let $\lambda_0=\max(0,-\lambda_{\min}(K))$. For
$\lambda>\lambda_0$, $\norm{(K+\lambda\Id)^{-1}\mathbf k}$ is continuous,
nonincreasing, and tends to zero. If its limit at $\lambda_0$ exceeds one,
choose a root with norm one. If the limit is at most one, the excluded
$\lambda=0$ case implies $\lambda_0>0$; the limiting matrix is singular
and $\mathbf k$ is orthogonal to its kernel. Add a kernel vector to its
minimum-norm solution to reach the unit sphere. The square term vanishes in
both cases, giving equality. The open ball has the same supremum as its
closure by continuity, but need not attain a boundary optimum.

In a $K$ eigenbasis, the scalar objective is
$\lambda+\sum_i k_i^2/(\kappa_i+\lambda)$; a zero denominator is admissible
only when the corresponding numerator vanishes. Unlike the formula in the
violating regime, this exact characterization retains a scalar minimization.
